\chapter{Mesurer le market making et l'exécution}\label{ch:rs:measuring-market-making-and-execution}

Un teneur de marché qui cote un lot de chaque côté d'une action à 100~\$ termine son année en perte de 945\,000~\$. Ses livres font apparaître une \omterm{def:rs:measuring-market-making-and-execution:pnl}{capture du spread} de $+\$480\,000$, une sélection adverse
de $-\$1.23$ million, des pertes d'inventaire de $-\$145\,000$ et des frais de $-\$56\,000$. L'année correspond à douze heures simulées de \texttt{firm.tape} ramenées à 252 jours
de 6.5 heures, et la décomposition est l'essentiel~: elle dit que les cotations rapportent leur demi-fourchette puis perdent plus de deux fois autant au profit des traders qui
les frappent en moins de vingt secondes~; que 42~\% des exécutions rencontrent des ordres informés, qui coûtent 1.85 tick par action, tandis que les exécutions non informées rapportent de l'argent~; et
qu'un filtre qui négocie 59~\% de moins réduit la perte de 61~\% sans rendre une seule exécution nettement meilleure. Ce chapitre mesure les exécutions d'un teneur de marché et
d'un algorithme d'exécution, avec \texttt{firm.markout}~: \omterm{def:rs:measuring-market-making-and-execution:curve}{courbes de mark-out}, statistiques d'exécution, une décomposition du P\&L qui se somme exactement, et \omterm{def:rs:measuring-market-making-and-execution:tca}{analyse des coûts de transaction} (TCA).

\section{Courbes de mark-out}

\begin{definition}[Courbe de mark-out]\label{def:rs:measuring-market-making-and-execution:curve}
La \emph{courbe de mark-out}\index{courbe de mark-out} d'un ensemble d'exécutions est la moyenne pondérée par les quantités, en fonction de l'horizon $h$, du mark-out de chaque exécution
(livre~2, chapitre~15), $s_i\,(m(t_i + h) - p_i)$ pour une exécution de sens $s_i$ au prix $p_i$ et à l'instant $t_i$, par rapport à un prix de référence $m$~: le mid, ou le
\omterm{def:rs:order-book-features:micro}{microprix} (chapitre~8).
\end{definition}

Le teneur de marché du chapitre évolue dans \texttt{firm.tape} comme substitut en direct de celui du chapitre~19~: un lot au meilleur cours acheteur et un lot au meilleur cours vendeur, dans la limite de cinq lots
d'inventaire, avec des latences qu'il tire lui-même (données de marché 0.02 seconde plus une loi exponentielle de moyenne 0.02, saisie des ordres 0.03 plus 0.02) et des frais de 0.05 tick par
action. Sur douze séances d'une heure, il affiche 18\,000 lots et 8\,144 sont exécutés. La contrepartie de chaque exécution est connue, car la bande enregistre si
l'ordre au marché qui a touché la cotation était informé. La \Cref{fig:rs:measuring-market-making-and-execution:curve} montre les mark-out~: 0.43 tick par action à l'exécution
par rapport au mid juste avant elle (la demi-fourchette, moins lorsque la fourchette valait un tick et que l'exécution intervenait à un niveau périmé), tombant à $-0.67$ à vingt secondes.
Les exécutions non informées conservent un mark-out positif, 0.18 à vingt secondes et 0.52 à cinq minutes~; les informées tombent à $-1.85$ et $-2.71$. Par rapport au \omterm{def:rs:order-book-features:micro}{microprix}
l'exécution ne rapporte que 0.19 tick à $h = 0$~: le \omterm{def:rs:order-book-features:micro}{microprix} penche déjà vers là où va le mid, et mesure la fourchette qui a réellement été captée.

L'horizon de la courbe compte autant que son niveau. Jusqu'à quelques secondes, le mark-out mêle fourchette et bruit~; au-delà d'un certain horizon, il cesse de bouger, parce que
l'information contenue dans la transaction a été absorbée. Pour le cotateur du chapitre, la courbe agrégée reste à moins de deux erreurs types de sa valeur à cinq minutes à partir de vingt
secondes~: les mark-out se stabilisent à vingt secondes, et c'est l'horizon auquel la sélection adverse se mesure. Brogaard, Hendershott et Riordan retrouvent la
même forme dans les données réelles~: les ordres apporteurs de liquidité des traders haute fréquence subissent une sélection adverse, et leurs transactions prédisent les variations de prix sur des horizons de
secondes.

\begin{omfigure}
\begin{tikzpicture}
\begin{semilogxaxis}[width=11.5cm,height=6cm,xlabel={horizon (secondes~; le premier point est l'exécution)},ylabel={mark-out (ticks par action)},
  tick label style={font=\small},label style={font=\small},xmin=0.04,xmax=400,ymin=-3,ymax=1,grid=major,grid style={gray!25},
  xtick={0.05,0.1,1,10,100},xticklabels={exéc.,0.1,1,10,100},legend columns=2,
  legend style={font=\footnotesize,at={(0.5,-0.3)},anchor=north,draw=none,fill=none,/tikz/every even column/.append style={column sep=8pt}},
  table/col sep=comma]
\addplot[name path=lo,draw=none,forget plot] table[x=h,y=all_lo]{figdata/research/23-measuring-market-making-and-execution/markout.csv};
\addplot[name path=hi,draw=none,forget plot] table[x=h,y=all_hi]{figdata/research/23-measuring-market-making-and-execution/markout.csv};
\addplot[omDef!20,forget plot] fill between[of=lo and hi];
\addplot[omDef,very thick,mark=*,mark size=1.2pt] table[x=h,y=all]{figdata/research/23-measuring-market-making-and-execution/markout.csv};
\addplot[omDef,thick,dashed] table[x=h,y=micro]{figdata/research/23-measuring-market-making-and-execution/markout.csv};
\addplot[omThm,thick,mark=triangle*,mark size=1.4pt] table[x=h,y=informed]{figdata/research/23-measuring-market-making-and-execution/markout.csv};
\addplot[black!60,thick,mark=square*,mark size=1pt] table[x=h,y=uninformed]{figdata/research/23-measuring-market-making-and-execution/markout.csv};
\draw[black!50,densely dotted] (axis cs:20,-3) -- (axis cs:20,1);
\legend{toutes les exécutions (mid),toutes les exécutions (microprix),contrepartie informée,contrepartie non informée}
\end{semilogxaxis}
\end{tikzpicture}

\omcaption{\omterm{def:rs:measuring-market-making-and-execution:curve}{Courbes de mark-out} des exécutions du cotateur au meilleur cours sur douze heures simulées, avec une bande de deux erreurs types pour l'ensemble des exécutions. La ligne pointillée marque vingt
secondes, où la courbe se stabilise. Données~: \texttt{rs\_markout.curves}.}\label{fig:rs:measuring-market-making-and-execution:curve}
\end{omfigure}

\section{Taux d'exécution et taux de touche}

\begin{definition}[Taux d'exécution, taux de touche]\label{def:rs:measuring-market-making-and-execution:fill}
Le \emph{taux d'exécution}\index{taux d'exécution} d'un ensemble d'ordres passifs est la quantité exécutée rapportée à la quantité affichée. Le \emph{taux de touche}\index{taux de touche}
d'un dealer répondant aux demandes de cotation (RFQ) est la part de ses cotations que les clients touchent.
\end{definition}

Un \omterm{def:rs:measuring-market-making-and-execution:fill}{taux d'exécution} indique quelle part de ce que le teneur de marché offrait a été prise, mais il ne dit pas à lui seul si c'est bon~: les cotations sont prises le plus volontiers lorsqu'elles
sont fausses. Le second cotateur du chapitre retire son cours acheteur chaque fois que la file acheteuse compte moins de 30~\% de la taille au meilleur cours, et son cours vendeur de manière symétrique,
car une file peu fournie annonce que le prix la traversera (chapitre~8)~:

\begin{center}
\small
\begin{tabular}{@{}lrr@{}}
\toprule
 douze heures simulées & cotateur au meilleur cours & avec le filtre de déséquilibre\\
\midrule
 lots affichés, lots exécutés & 18\,000, 8\,144 & 16\,391, 3\,309\\
 \omterm{def:rs:measuring-market-making-and-execution:fill}{taux d'exécution} & 45~\% & 20~\%\\
 part des exécutions face aux ordres informés & 42~\% & 40~\%\\
 mark-out à l'exécution, par rapport au mid & 0.43 & 0.52\\
 \quad par rapport au \omterm{def:rs:order-book-features:micro}{microprix} & 0.19 & 0.44\\
 mark-out à vingt secondes & $-0.67$ & $-0.53$\\
 P\&L par lot exécuté & $-\$0.85$ & $-\$0.81$\\
\bottomrule
\end{tabular}
\end{center}

Le filtre exécute moins de moitié aussi souvent, et ses exécutions ont lieu lorsque le \omterm{def:rs:order-book-features:micro}{microprix} se trouve près du mid, de sorte qu'il capte davantage la fourchette (0.44 tick par rapport
au \omterm{def:rs:order-book-features:micro}{microprix} plutôt que 0.19)~; mais les traders qui le touchent sont aussi souvent informés qu'auparavant, et par lot il perd presque autant. Le \omterm{def:rs:measuring-market-making-and-execution:fill}{taux de touche} d'un dealer
présente la même ambiguïté sur une plateforme de demandes de cotation~: un dealer qui remporte chaque demande cote trop bien.

\section{Le P\&L du teneur de marché, décomposé}

\begin{definition}[Capture du spread, coût de sélection adverse, P\&L d'inventaire]\label{def:rs:measuring-market-making-and-execution:pnl}
Pour une exécution de sens $s$, de quantité $q$ et de prix $p$ à l'instant $t$, avec un prix de référence $m$, un horizon $H$ et la fin de la période $T$~: la \emph{capture du
spread}\index{capture du spread} est $sq\,(m(t) - p)$~; le \emph{coût de sélection adverse}\index{coût de sélection adverse} est $sq\,(m(t + H) - m(t))$~; le \emph{P\&L
d'inventaire}\index{P\&L d'inventaire} est $sq\,(m(T) - m(t + H))$. Sommés sur les exécutions, frais inclus, les trois termes font le P\&L des exécutions, la position finale étant valorisée à
$m(T)$.
\end{definition}

L'identité est exacte~: $m(T) - p = (m(t) - p) + (m(t+H) - m(t)) + (m(T) - m(t+H))$ pour chaque exécution (le \cref{lst:rs:measuring-market-making-and-execution:decompose}).
Ce que l'identité laisse à l'appréciation est l'horizon $H$~: trop court, et la sélection adverse qui se déploie encore après cet horizon est enregistrée comme inventaire~; trop long, et le risque d'inventaire
que le teneur de marché a choisi de porter est enregistré comme sélection adverse. L'horizon de stabilisation de la \omterm{def:rs:measuring-market-making-and-execution:curve}{courbe de mark-out} est le choix naturel, ici vingt secondes.
La décomposition de Menkveld d'un teneur de marché haute fréquence réel sur Chi-X et Euronext a la même structure~: un gain de \texteuro1.55 par transaction sur la fourchette, nette
des frais, une perte de positionnement de \texteuro0.68 (un gain sur les positions tenues moins de cinq secondes, une perte sur les plus longues), un gain brut de \texteuro0.88.

\begin{omfigure}
\begin{tikzpicture}
\begin{axis}[width=11cm,height=5.8cm,ybar,bar width=12pt,ylabel={dollars, douze heures},symbolic x coords={spread,adverse,inventory,fees,total},
  xtick=data,xticklabels={capture du\\spread,sélection\\adverse,inventaire,frais,total},xticklabel style={align=center},
  tick label style={font=\small},label style={font=\small},enlarge x limits=0.15,
  ymin=-10000,ymax=5000,grid=major,grid style={gray!25},legend columns=2,
  yticklabel style={/pgf/number format/fixed,/pgf/number format/1000 sep={\,}},scaled y ticks=false,
  legend style={font=\footnotesize,at={(0.5,-0.3)},anchor=north,draw=none,fill=none,/tikz/every even column/.append style={column sep=8pt}},
  table/col sep=comma]
\addplot[fill=omDef!55,draw=omDef] table[x=part,y=plain]{figdata/research/23-measuring-market-making-and-execution/decomp.csv};
\addplot[fill=gray!45,draw=gray] table[x=part,y=pull]{figdata/research/23-measuring-market-making-and-execution/decomp.csv};
\legend{cotateur au meilleur cours,avec le filtre de déséquilibre}
\end{axis}
\end{tikzpicture}

\omcaption{Le P\&L des deux cotateurs sur douze heures simulées, décomposé sur un horizon de vingt secondes. Données~:
\texttt{rs\_markout.decomposition}.}\label{fig:rs:measuring-market-making-and-execution:decomp}
\end{omfigure}

Sur les douze heures, le cotateur au meilleur cours capture 3\,519~\$ de spread, perd 8\,977~\$ au titre de la sélection adverse dans les vingt secondes, 1\,060~\$ sur l'inventaire qui suit, et
paie 407~\$ de frais~: une perte de 6\,925~\$ (la \cref{fig:rs:measuring-market-making-and-execution:decomp}). Par action, le spread rapporte 0.43 tick et la sélection adverse
coûte 1.10, soit le mark-out de $-0.67$ à vingt secondes. Le cotateur filtré capture 1\,734~\$, perd 3\,483~\$ et 771~\$, paie 165~\$~: $-\$2\,685$. Aucun des deux
n'est une activité viable~; la décomposition indique ce qu'il faudrait modifier (moins de contreparties informées, ou un spread plus large) et que l'inventaire n'est pas le problème.

\section{Mesurer l'exécution}

\begin{definition}[Coût de délai, coût d'opportunité, slippage par rapport au VWAP]\label{def:rs:measuring-market-making-and-execution:costs}
Pour un ordre parent décidé au prix $P_d$, lancé au \omterm{def:rs:simulation-versus-live:calibration}{prix d'arrivée} $P_a$, exécuté à hauteur de $Q$ sur sa cible $X$ au prix moyen $\bar P$, avec le prix $P_e$
à sa fin~: le \emph{coût de délai}\index{coût de délai} vaut $sQ(P_a - P_d)$, le coût d'exécution $sQ(\bar P - P_a)$, le \emph{coût d'opportunité}\index{coût d'opportunité}
$s(X - Q)(P_e - P_d)$~; avec les frais, ils s'additionnent pour former l'\omterm{def:rs:simulation-versus-live:calibration}{implementation shortfall} (chapitre~19). Le \emph{slippage par rapport au VWAP}\index{slippage par rapport au VWAP} vaut $s(\bar P - \mathrm{VWAP})$,
soit le prix moyen comparé au \omterm{def:rs:market-data-for-research:bar}{prix moyen pondéré par les volumes} du marché sur la durée de vie de l'ordre.
\end{definition}

L'ordre parent du chapitre achète 300 lots~: décidé à la minute 10, lancé à la minute 11, à exécuter avant la minute 30. Toutes les vingt secondes, l'exécuteur affiche les lots de l'étape suivante
au meilleur cours acheteur et envoie un ordre au marché pour tout retard de plus de huit lots sur une trajectoire linéaire. Dans une session (graine 61), il exécute 28\,700
actions, dont 43~\% passivement en moyenne sur les sessions, et son rapport affiche~:

\begin{center}
\small
\begin{tabular}{@{}lr@{}}
\toprule
\omterm{def:rs:measuring-market-making-and-execution:tca}{analyse des coûts de transaction}, un ordre & ticks par action\\
\midrule
\omterm{def:rs:measuring-market-making-and-execution:costs}{coût de délai} & 0.00\\
coût d'exécution par rapport à l'arrivée & $-7.83$\\
\omterm{def:rs:measuring-market-making-and-execution:costs}{coût d'opportunité} (1\,300 actions non exécutées) & $-0.91$\\
frais & 0.05\\
\omterm{def:rs:simulation-versus-live:calibration}{implementation shortfall} & $-8.69$\\
\omterm{def:rs:measuring-market-making-and-execution:costs}{slippage par rapport au VWAP} & $-1.10$\\
\bottomrule
\end{tabular}
\end{center}

Un shortfall négatif de 8.69 ticks (8.7 points de base à 100~\$) ressemble à une réussite de l'exécution. C'est le marché~: le mid a baissé pendant ces dix-neuf minutes, et
un acheteur qui arrive avant une baisse paraît brillant face au \omterm{def:rs:simulation-versus-live:calibration}{prix d'arrivée}. Sur les douze sessions, le shortfall moyen vaut $-2.00$ ticks par action, avec une erreur type de
2.73~; sa part d'exécution, $-4.68$ (2.09), se décompose en la dérive du mid entre l'arrivée et chaque exécution, $-4.78$ (2.07), et le prix payé par rapport au mid juste
avant l'exécution, $+0.10$ (0.03). Seule la seconde relève de l'algorithme, et seule elle est mesurée avec précision. Même le \omterm{def:rs:measuring-market-making-and-execution:costs}{coût de délai}, $+2.83$ (0.94), relève du hasard~: les
mêmes douze sessions, menées sans l'ordre, montent de 2.92 ticks sur cette minute. Et le simulateur permet ce qu'aucun desk ne peut faire~: rejouer la même session sans
l'ordre. Le mouvement du mid entre la minute 11 et la minute 30 vaut $-9.71$ ticks avec l'ordre et $-9.96$ sans lui, soit un impact de $+0.25$ (0.24). Sur ce marché, le prix
efficient ne réagit pas à l'ordre, si bien que son impact durable est nul par construction~; celui d'un marché réel ne le serait pas (chapitre~27).

\begin{omfigure}
\begin{tikzpicture}
\begin{axis}[width=9.5cm,height=6cm,xlabel={mouvement du mid, minute 11 à 30, sans l'ordre (ticks)},ylabel={ticks par action},tick label style={font=\small},
  label style={font=\small},xmin=-60,xmax=20,ymin=-25,ymax=12,grid=major,grid style={gray!25},legend pos=north west,legend style={font=\footnotesize},
  table/col sep=comma]
\addplot[only marks,mark=*,mark size=2pt,omDef] table[x=move_without,y=execution]{figdata/research/23-measuring-market-making-and-execution/tca.csv};
\addplot[only marks,mark=square*,mark size=1.8pt,omThm] table[x=move_without,y=paid]{figdata/research/23-measuring-market-making-and-execution/tca.csv};
\legend{coût d'exécution par rapport à l'arrivée,prix payé par rapport au mid}
\end{axis}
\end{tikzpicture}

\omcaption{Douze ordres parents~: le coût d'exécution mesuré par rapport au \omterm{def:rs:simulation-versus-live:calibration}{prix d'arrivée} suit le mouvement propre du marché sur la durée de vie de l'ordre (mesuré dans la même
session menée sans l'ordre)~; le prix payé par rapport au mid à chaque exécution ne le suit pas. Données~: \texttt{rs\_markout.tca\_all}.}\label{fig:rs:measuring-market-making-and-execution:tca}
\end{omfigure}

\section{Analyse des coûts de transaction}

\begin{definition}[Analyse des coûts de transaction]\label{def:rs:measuring-market-making-and-execution:tca}
L'\emph{analyse des coûts de transaction}\index{analyse des coûts de transaction} (TCA) est la mesure des coûts des ordres exécutés par rapport à des indices de référence (les prix de décision et
d'arrivée, le VWAP, le cours de clôture), décomposée en parts imputables aux décisions (le délai, le choix de laisser une quantité non exécutée) et à l'exécution (le prix payé
par rapport au marché à chaque exécution), et agrégée sur de nombreux ordres pour comparer les courtiers, les algorithmes et les traders.
\end{definition}

La leçon des douze ordres est que la TCA est un exercice statistique au bruit important. Un shortfall par rapport au \omterm{def:rs:simulation-versus-live:calibration}{prix d'arrivée} a la variance du mouvement du marché
sur la durée de vie de l'ordre, qui écrase le coût mesuré~; les moyennes sur de nombreux ordres, conditionnées au mouvement du marché (en régressant le shortfall sur le
rendement contemporain de l'action ou de son secteur), et les mesures de coût ancrées à chaque exécution (le prix payé par rapport au mid, le mark-out propre à l'exécution) sont ce qui
la rend informative. Le \omterm{def:rs:measuring-market-making-and-execution:costs}{slippage par rapport au VWAP} est moins bruité (0.40 tick avec une erreur type de 0.61 ici), parce que l'indice de référence suit le marché, et il peut être manipulé
pour la même raison~: un algorithme qui échange au prorata des volumes égale le VWAP quel qu'en soit le coût. Un rapport de TCA doit indiquer l'indice de référence, la décomposition,
l'erreur type et le nombre d'ordres, et comparer les algorithmes par expérience randomisée (chapitre~21).

\section{Tutoriel~: l'année plate}\label{tut:rs:measuring-market-making-and-execution:tut}

\begin{tutorial}
\textbf{Objectif.} Faire tourner deux teneurs de marché et un exécuteur dans \texttt{firm.tape}, tracer les \omterm{def:rs:measuring-market-making-and-execution:curve}{courbes de mark-out} par contrepartie, décomposer le P\&L, et rédiger
la TCA d'un ordre parent. \textbf{État final~:} les deux tableaux et
\cref{fig:rs:measuring-market-making-and-execution:curve,fig:rs:measuring-market-making-and-execution:decomp,fig:rs:measuring-market-making-and-execution:tca}.
\begin{tutsteps}
\item \textbf{Mark-outs}~: le prix de référence en vigueur à chaque instant, et chaque exécution comparée à ce prix à chaque horizon.
\omcode{code/firm/markout/firm_markout.py}{38}{45}{Mark-outs par rapport à une trajectoire de prix de référence.}\label{lst:rs:measuring-market-making-and-execution:markouts}
\item \textbf{La décomposition} qui s'additionne.
\omcode{code/firm/markout/firm_markout.py}{77}{90}{Capture du spread, sélection adverse, inventaire, frais.}\label{lst:rs:measuring-market-making-and-execution:decompose}
\item \textbf{Le shortfall} d'un ordre parent.
\omcode{code/firm/markout/firm_markout.py}{93}{103}{Implementation shortfall en quatre parts.}\label{lst:rs:measuring-market-making-and-execution:shortfall}
\item \textbf{Exécuter} \texttt{rs\_markout.curves}, \texttt{decomposition}, \texttt{settle}, \texttt{execution}, \texttt{tca\_all} et \texttt{fig\_markout.py}.
\end{tutsteps}
\textbf{Pour aller plus loin.} Coter un tick derrière le meilleur cours et constater que le \omterm{def:rs:measuring-market-making-and-execution:fill}{taux d'exécution} baisse et que le mark-out par lot augmente~; donner à l'exécuteur un signal (le
\omterm{def:rs:order-book-features:micro}{microprix}) et mesurer si son prix payé par rapport au mid s'améliore.
\end{tutorial}

\section{Construction~: le module de mesure}\label{bld:rs:measuring-market-making-and-execution:markout}

\begin{build}
\bfield{Rôle} Chaque exécution de la firme, passive ou agressive, est mesurée de la même façon~: mark-outs aux horizons standard par rapport au mid et au \omterm{def:rs:order-book-features:micro}{microprix},
regroupés par contrepartie, plateforme et stratégie~; P\&L de teneur de marché décomposé~; chaque ordre parent analysé.
\bfield{Interface} \texttt{microprice}, \texttt{ref\_at(ref\_t, ref\_px, t)}, \texttt{markouts(t, side, px, ref\_t, ref\_px, horizons)},
\texttt{curve(M, qty, groups)}, \texttt{settle\_horizon(mean, horizons, tol)}, \texttt{fill\_rate}, \texttt{hit\_ratio}, \texttt{mm\_decompose(t, side, px, qty, ref\_t, ref\_px, H, fee,
t\_end)}, \texttt{shortfall(side, target, decision\_px, arrival\_px, fills, end\_px, fee)}, \texttt{vwap\_slippage}, \texttt{tca\_report}.
\bfield{Règles} La référence est celle en vigueur avant l'exécution~; la décomposition s'additionne au P\&L marqué au centime près~; les horizons sont fixés à l'échelle de la firme~;
chaque moyenne est rapportée avec son erreur type et son effectif.
\bfield{Tests d'acceptation} \texttt{code/firm/markout/tests/}~: références et mark-outs calculés à la main~; une courbe pondérée~; l'horizon de règlement~; la décomposition égale au
P\&L d'exécutions aléatoires avec la position marquée à la fin~; le shortfall égal au gain du portefeuille papier moins celui du portefeuille réel~; \omterm{def:rs:measuring-market-making-and-execution:costs}{slippage par rapport au VWAP}, \omterm{def:rs:measuring-market-making-and-execution:fill}{taux d'exécution} et
\omterm{def:rs:measuring-market-making-and-execution:fill}{taux de touche} calculés à la main.
\bfield{Pour aller plus loin} Mark-outs par rapport à un modèle de fair value~; TCA ajustée par régression (shortfall sur le mouvement du marché)~; cartes de notation des contreparties pour un dealer.
\end{build}

\begin{omsources}
\item A.~J.~Menkveld, ``High frequency trading and the new market makers'', \emph{Journal of Financial Markets} 16(4), 2013 (Tinbergen Institute discussion
paper 11-076, 2011).
\item J.~Brogaard, T.~Hendershott and R.~Riordan, ``High-frequency trading and price discovery'', \emph{Review of Financial Studies} 27(8), 2014.
\item A.~F.~Perold, ``The implementation shortfall: paper versus reality'', \emph{Journal of Portfolio Management} 14(3), 1988.
\end{omsources}

\section{Exercices}

\begin{exercise}[$\star$]\label{exo:rs:measuring-market-making-and-execution:1}
Un achat exécuté à 99.99 alors que le milieu était à 100.00~; vingt secondes plus tard le milieu vaut 99.97, en fin de journée 100.02. Décomposer le P\&L par action de l'exécution en
\omterm{def:rs:measuring-market-making-and-execution:pnl}{capture du spread}, sélection adverse et inventaire (en ticks de 0.01).
\end{exercise}

\begin{exercise}[$\star$]\label{exo:rs:measuring-market-making-and-execution:2}
Le teneur de cotations au meilleur niveau a affiché 18\,000 lots et 8\,144 ont été exécutés. Quel est son \omterm{def:rs:measuring-market-making-and-execution:fill}{taux d'exécution}~? Pourquoi un \omterm{def:rs:measuring-market-making-and-execution:fill}{taux d'exécution} plus élevé n'est-il pas meilleur~?
\end{exercise}

\begin{exercise}[$\star$]\label{exo:rs:measuring-market-making-and-execution:3}
Le spread rapporte 0.43 tick par action et la sélection adverse à vingt secondes coûte 1.10. Quel est le mark-out à vingt secondes~?
\end{exercise}

\begin{exercise}[$\star\star$]\label{exo:rs:measuring-market-making-and-execution:4}
Pourquoi le choix de l'horizon $H$ ne modifie-t-il pas le total de la décomposition~? Quelles composantes déplace-t-il, et dans quel sens lorsque $H$ est trop court~?
\end{exercise}

\begin{exercise}[$\star\star$]\label{exo:rs:measuring-market-making-and-execution:5}
Un ordre d'achat de 100 actions est décidé à 50.00, transmis à 50.20, exécuté à hauteur de 80 actions à un prix moyen de 50.40~; le cours vaut 51.00 à sa fin~; les frais sont de 0.01 par action.
Calculer les \omterm{def:rs:measuring-market-making-and-execution:costs}{coûts de délai}, d'exécution et d'opportunité, ainsi que le shortfall.
\end{exercise}

\begin{exercise}[$\star\star$]\label{exo:rs:measuring-market-making-and-execution:6}
Le filtre de déséquilibre a réduit la perte de 6\,925~\$ à 2\,685~\$. Quelle part de cette réduction tient à moins d'exécutions et quelle part à de meilleures exécutions~?
\end{exercise}

\begin{exercise}[$\star\star\star$]\label{exo:rs:measuring-market-making-and-execution:7}
\emph{Programmation.} Étendre à une année de 252 jours de 6.5 heures la décomposition sur douze heures du teneur de cotations au meilleur niveau faite dans le chapitre, et préciser ce que cette extension suppose.
\end{exercise}

\begin{exercise}[$\star\star\star$]\label{exo:rs:measuring-market-making-and-execution:8}
\emph{Trouver la faille.} «~Notre nouvel algorithme d'exécution a battu le \omterm{def:rs:simulation-versus-live:calibration}{prix d'arrivée} de 8.7 points de base sur son premier ordre, il est donc le meilleur algorithme que nous ayons.~»
\end{exercise}

\section{Problème~: l'année plate}

\begin{problem}[{Problème du week-end --- les livres du teneur de marché}]\label{pb:rs:measuring-market-making-and-execution:1}
Les deux teneurs de cotations du chapitre et son exécutant, sur douze heures simulées de \texttt{firm.tape}.

\medskip\noindent\textbf{Partie I --- Mark-outs.}
\begin{enumerate}
\item Quel est le mark-out à l'exécution du teneur de cotations au meilleur niveau, par rapport au milieu et par rapport au \omterm{def:rs:order-book-features:micro}{microprix}, et pourquoi diffèrent-ils~?
\item Quels sont les mark-outs informés et non informés à vingt secondes et à cinq minutes~?
\item Quelle part des exécutions a rencontré des ordres informés~?
\item À quel horizon la courbe se stabilise-t-elle, et selon quelle règle~?
\end{enumerate}

\medskip\noindent\textbf{Partie II --- Exécutions.}
\begin{enumerate}[resume]
\item Donner les lots affichés et les lots exécutés par chacun des deux teneurs, ainsi que leurs \omterm{def:rs:measuring-market-making-and-execution:fill}{taux d'exécution}.
\item Que change le filtre de déséquilibre dans les mark-outs, et que ne change-t-il pas~?
\item Quel est le P\&L par lot de chaque teneur~?
\end{enumerate}

\medskip\noindent\textbf{Partie III --- La décomposition.}
\begin{enumerate}[resume]
\item Décomposer les douze heures du teneur de cotations au meilleur niveau à vingt secondes.
\item Décomposer celles du teneur filtré.
\item Comment les composantes spread et sélection adverse par action se rapportent-elles à la \omterm{def:rs:measuring-market-making-and-execution:curve}{courbe de mark-out}~?
\item À quoi ressemble la décomposition de Menkveld pour un teneur de marché réel~?
\end{enumerate}

\medskip\noindent\textbf{Partie IV --- Exécution et verdict.}
\begin{enumerate}[resume]
\item Donner le rapport d'\omterm{def:rs:measuring-market-making-and-execution:tca}{analyse des coûts de transaction} (TCA) d'un ordre pour la graine 61.
\item Moyenner le shortfall et ses composantes sur les douze ordres, avec leurs erreurs types.
\item Quelle part relève du travail de l'algorithme~?
\item Quel est l'impact de l'ordre sur le milieu, mesuré par le contrefactuel, et pourquoi est-il faible ici~?
\item Énoncer le \emph{résultat nommé}~: l'année du teneur de cotations au meilleur niveau décomposée en \omterm{def:rs:measuring-market-making-and-execution:pnl}{capture du spread}, sélection adverse, inventaire et frais, et l'horizon auquel ses
mark-outs se stabilisent.
\item Que devrait changer le teneur de marché pour atteindre le seuil de rentabilité~?
\item Comment un desk doit-il comparer deux algorithmes d'exécution~?
\item Que doit énoncer chaque rapport de TCA~?
\item En une phrase~: que mesure une \omterm{def:rs:measuring-market-making-and-execution:curve}{courbe de mark-out}~?
\end{enumerate}
\end{problem}

\section{Questions d'entretien}

\begin{interviewq}[$\star$ \iqroles{trader}]\label{iq:rs:measuring-market-making-and-execution:1}
Qu'est-ce qu'un mark-out, et à quel horizon le calculeriez-vous~?
\end{interviewq}

\begin{interviewq}[$\star\star$ \iqroles{trader, chercheur}]\label{iq:rs:measuring-market-making-and-execution:2}
Décomposeriez-vous le P\&L d'un teneur de marché~? Comment choisissez-vous la répartition entre sélection adverse et inventaire~?
\end{interviewq}

\begin{interviewq}[$\star\star$ \iqroles{chercheur}]\label{iq:rs:measuring-market-making-and-execution:3}
Un teneur de marché à haute fréquence gagne \texteuro1.55 par transaction sur le spread et perd \texteuro0.68 sur ses positions. Qu'est-ce que cela vous dit de son activité~?
\end{interviewq}

\begin{interviewq}[$\star\star$ \iqroles{trader, banque}]\label{iq:rs:measuring-market-making-and-execution:4}
Votre \omterm{def:rs:measuring-market-making-and-execution:fill}{taux de touche} sur une plateforme de demandes de cotation a doublé le mois dernier. Est-ce une bonne nouvelle~?
\end{interviewq}

\begin{interviewq}[$\star\star$ \iqroles{trader, chercheur}]\label{iq:rs:measuring-market-making-and-execution:5}
Comment savez-vous si un algorithme d'exécution est meilleur qu'un autre, sachant que le shortfall de chaque ordre est dominé par le mouvement du marché~?
\end{interviewq}

\begin{interviewq}[$\star\star\star$ \iqroles{chercheur, développeur}]\label{iq:rs:measuring-market-making-and-execution:6}
Concevez le service de mesure des exécutions de la firme~: entrées, prix de référence, horizons, regroupements, et contrôles vérifiant que ses chiffres se recoupent.
\end{interviewq}
